\scanpage{58}
application $L\mathrel{\overset{g}{-}}K\times K'$ est $\mathcal M_0$-modérée ssi $L\mathrel{\overset{f}{-}}K$ et $L\mathrel{\overset{f'}{-}}K'$ le sont. En effet, \deleted{si $g$ est}
\diagram{\begin{tikzcd}[ampersand replacement=\&]
K\times K'\arrow[r,dash,"{\mathrm{pr}_1}"]\arrow[dr,dash,"{\mathrm{pr}_2}"']\&K\\
\&K'
\end{tikzcd}}
sont modérés car le graphe de $K\times K'\mathrel{\overset{\mathrm{pr}_1}{-}}K$ disons est
\[
\deleted{\Delta_K\times K'\subset I^{\unsure{2n+n'}}}\quad
(K\times K'\times K)\cap
\mathrm{Im}(I^{n+n'}\longrightarrow I^{n+n'+n})
\]
\[
(x,x')\longmapsto(x,x',x)
\]
qui est $\in M_{2n+n'}$ par M'1,2,3. Donc $g$ modéré $\Rightarrow f,f'$ modérés.
Inv. On a
\[
\Gamma_g\subset I^{m+n+n'}
=p^{-1}(\Gamma_f)\cap q^{-1}(\unsure{\Gamma_g|\Gamma_{f'}}),
\]
\[
p:I^{m+n+n'}-I^{m+n},\qquad q:I^{m+n+n'}-I^{m+n'},
\]
et appliquer M'1,2. Le foncteur \deleted{$\unsure{\mathcal M_0\to\text{esp.}}$} \inserted{$(K,n)\longmapsto K$} commute aux produits finis.

On voit de façon analogue que si
\[
f,g:(K,n)\rightrightarrows(K',n'),
\]
alors $\mathrm{Ker}(f,g)$ existe dans $\mathcal M_0$ et $\mathcal M_0-(\text{esp cpts})$ y commute, car considérons
\[
\mathrm{Ker}(f,g)=r\left(p^{-1}(\Gamma_f)\cap q^{-1}(\Gamma_g)
\inserted{\cap\delta(I^{n+n'})}\right)
\quad\text{est }\in M_{n+n'+n'},
\]
où $p,q:I^{n+n'+n'}-I^{n+n'}$ sont les 2 proj.
\[
\delta:I^{n+n'}-I^{n+n'+n'},\qquad
r:I^{n+n'+n'}\longrightarrow I^{\unsure{n}},
\]

\diagram{\begin{tikzcd}[ampersand replacement=\&]
K\arrow[r,"{f,g}"]\&K'\times K'\subset I^{n'+n'}\\
\&I^{n'}\arrow[u,"{\delta}"]
\end{tikzcd}}

\margin{$x,x'_1,x'_2$\\
$x'_1=f(x)$\\
$x'_2=f'(x)$\\
$x'_1=x'_2$}

Ainsi on a M1), M2), et on a aussi M3) grâce à M'3). On a l'objet $(I,1)=\underline I$, \inserted{de $\mathcal M_0$,} et l'axiome M5 par construction.

Alors les conditions supplémentaires b) avant M4), et M4), équivalent à :

M'4) $0^\circ$) \deleted{\unsure{Pour}} $q\in\mathbb Q$, $|q|\leq1$, l'ens
$\{(x,y)\mid y=qx\}\subset I\times I$ est $\in M_2$

$1^\circ$) \deleted{$\{1\}\subset I$} $\{1\}\subset I$ est $\in M_1$

$2^\circ$) l'ens $\{x,y,z\mid z=\dfrac{x+y}{2}\}\subset I^3$ est $\in M_3$

$3^\circ$) l'ens $\{x,y,z\mid z=xy\}$ $(I^3)$ \text{---}.

\margin{59}
